\section{Boundary kernel calculations}
\label{app:boundary-methods}

The scalar calculation uses the same raised-cosine probability law and
normalization as section~\ref{sec:lk}.  At fixed sampled momenta, averaging
over the uniform phases gives
\begin{equation}
\widehat K_N(a,b)=\frac{2L_w}{N}\sum_{j=1}^N
\cos\bigl(k_j\cdot(a-b)\bigr).
\end{equation}
The reflected conditional covariance is the four-term expression
\begin{equation}
\widehat K_N^\sigma(z,w)=\tfrac12\bigl[
\widehat K_N(z,w)+\sigma\widehat K_N(z,\bar w)
+\sigma\widehat K_N(\bar z,w)+\widehat K_N(\bar z,\bar w)\bigr].
\label{eq:boundary-four-term-covariance}
\end{equation}
On the diagonal, its difference between heights reduces exactly to
\begin{equation}
\widehat V_\sigma(y;y_0)=\frac{2\sigma L_w}{N}\sum_{j=1}^N
\bigl[\cos(2k_j^y y)-\cos(2k_j^y y_0)\bigr].
\label{eq:boundary-feature-estimator}
\end{equation}
Thus the calculation samples the frequency-dependent conditional
covariance, rather than estimating products of fields from a finite
number of phase realizations.  Radial momenta are drawn by rejection
sampling from $W_\tau(u)$ in log radius and angles are uniform.  We use
$N=16\,384$, 32 independent batches, seed 20261107, 17 logarithmically
spaced heights $0.1\le y\le1.6$, and $y_0=0.4$.  Each batch uses common
momenta at all heights.  A straight-line fit with an intercept in
$\log(y/y_0)$ is performed in each batch; the standard error quoted in
section~\ref{sec:boundary} is the sample standard deviation of those 32
slopes divided by $\sqrt{32}$.  Correlations among the heights are thereby
retained.

The regulated reference is
\begin{equation}
V_\sigma^{\rm reg}(y;y_0)=2\sigma\int_{\Lambda_{\rm IR}}^{\Lambda_{\rm UV}}
\frac{dk}{k}\,W_\tau(u(k))
\bigl[J_0(2ky)-J_0(2ky_0)\bigr].
\end{equation}
We split the quadrature at the taper joins and into intervals resolving
the Bessel oscillations.  Tightening its tolerance and halving the maximum
interval length checks the numerical reference.  The largest deviation of
the regulated reference from $-2\log(y/y_0)$ is $2.8\times10^{-5}$
over these heights.

The fermion calculation uses the deterministic truncation already defined
in section~\ref{sec:susy},
\begin{equation}
C_L^>(u,v)=\frac{1-(v/u)^L}{u-v},\qquad |v|<|u|.
\end{equation}
For the mixed entry, $u=x+i\varepsilon$ and $v=-i\varepsilon$, so
the relative truncation error is $|v/u|^L$.  We take $L=32$,
$\varepsilon=0.05$, and 60 logarithmically spaced points
$0.1\le x\le3$.  The largest error is bounded by
$(1/\sqrt5)^{32}=6.5536\times10^{-12}$.
The exchange of the two ordered arguments includes the fermion minus
sign.  Reflected evaluations use the same mode features; no independent
right-sector draw is introduced.
